\documentclass[tikz,border=8pt]{standalone}
\usepackage{fontspec}
\usepackage{xeCJK}
\IfFontExistsTF{Segoe UI}{\setmainfont{Segoe UI}}{\setmainfont{Arial}}
\IfFontExistsTF{Microsoft JhengHei}{\setCJKmainfont{Microsoft JhengHei}}{\IfFontExistsTF{Noto Sans CJK TC}{\setCJKmainfont{Noto Sans CJK TC}}{\setCJKmainfont{Noto Sans TC}}}
\IfFontExistsTF{Microsoft JhengHei}{\setCJKsansfont{Microsoft JhengHei}}{\IfFontExistsTF{Noto Sans CJK TC}{\setCJKsansfont{Noto Sans CJK TC}}{\setCJKsansfont{Noto Sans TC}}}
\xeCJKsetup{CJKglue={\hskip 0.045em plus 0.015em minus 0.005em}, CJKecglue={\hskip 0.08em plus 0.02em}}
\IfFileExists{../../../FontAwesome.otf}{\newfontfamily\FA[Path=../../../]{FontAwesome.otf}}{\newfontfamily\FA{Segoe UI Symbol}}
\newcommand{\faIcon}[1]{{\FA\symbol{#1}}}
\usetikzlibrary{arrows.meta,positioning}
\definecolor{navy}{HTML}{0F172A}
\definecolor{muted}{HTML}{334155}
\definecolor{paper}{HTML}{FFFDF8}
\definecolor{blue}{HTML}{2563EB}
\definecolor{bluebg}{HTML}{DBEAFE}
\definecolor{green}{HTML}{00856A}
\definecolor{greenbg}{HTML}{DDFCEB}
\definecolor{gold}{HTML}{D97706}
\definecolor{goldbg}{HTML}{FEF3C7}
\definecolor{red}{HTML}{DC2626}
\definecolor{redbg}{HTML}{FEE2E2}
\definecolor{line}{HTML}{CBD5E1}
\begin{document}
\begin{tikzpicture}[x=1cm,y=1cm,>=Latex]
\path[fill=paper] (0,0) rectangle (16,9.6);
\node[font=\bfseries\fontsize{18}{22}\selectfont,text=navy,align=center,text width=14.4cm] at (8,8.82) {把主張放上壓力測試台};
\node[font=\fontsize{10.5}{13}\selectfont,text=muted,align=center,text width=13.8cm] at (8,8.12) {AI 不是幫我們把話說順，而是先逼我們把話說清楚};

\tikzset{
  box/.style={rounded corners=8pt,line width=1pt,align=center,minimum height=1.28cm,text width=2.55cm,font=\sffamily\fontsize{9.2}{11.4}\selectfont,text=navy},
  note/.style={rounded corners=7pt,draw=line,line width=.9pt,fill=white,align=center,text width=4.2cm,font=\sffamily\fontsize{8.4}{10.5}\selectfont,text=muted}
}
\node[box,draw=blue,fill=bluebg] (claim) at (2.25,5.1) {\bfseries 原始主張\\[-1pt]\normalfont 兩三句就好};
\node[box,draw=red,fill=redbg] (attack) at (5.65,5.1) {\bfseries AI 反方\\[-1pt]\normalfont 只攻最脆弱處};
\node[box,draw=gold,fill=goldbg] (evidence) at (9.05,5.1) {\bfseries 查資料\\[-1pt]\normalfont 分清證據與猜測};
\node[box,draw=green,fill=greenbg] (revision) at (12.45,5.1) {\bfseries 修正版本\\[-1pt]\normalfont 變窄、變準};
\draw[->,draw=navy,line width=1.1pt] (claim) -- (attack);
\draw[->,draw=navy,line width=1.1pt] (attack) -- (evidence);
\draw[->,draw=navy,line width=1.1pt] (evidence) -- (revision);

\node[note] at (3.95,3.22) {不要讓 AI 直接重寫；先讓學生交出會被追問的句子。};
\node[note] at (10.7,3.22) {評分看修正理由，不只看最後答案漂不漂亮。};
\draw[draw=line,line width=.9pt] (1.2,1.55)--(14.8,1.55);
\node[font=\fontsize{10.2}{12.8}\selectfont,text=navy,align=center,text width=13.8cm] at (8,.94) {好的反方不需要話多；一句打中弱點，就足夠讓我們回去改。};
\end{tikzpicture}
\end{document}
